\chapter{CRUX of RTL Files and Status}

\begingroup
\renewcommand{\arraystretch}{1.3}
\begin{longtable}{@{}p{0.22\textwidth}p{0.5\textwidth}p{0.22\textwidth}@{}}
\toprule
\textbf{Chapter} & \textbf{RTL Modules / Files} & \textbf{Status} \\
\midrule
\endhead

Chapter 9. Execution Stage &
\rtl{exe\_stage.sv}, \rtl{alu.sv}, \rtl{alu\_add.sv}, \rtl{alu\_logic.sv}, \rtl{alu\_shift.sv}, \rtl{alu\_cmp.sv}, \rtl{branch\_unit.sv}, \rtl{mul\_unit.sv}, \rtl{div\_unit.sv}, \rtl{mem\_unit.sv}, \rtl{vagu.sv}, \rtl{load\_store\_queue.sv}, \rtl{store\_buffer.sv}, \rtl{score\_board\_scalar.sv}, \rtl{score\_board\_simd.sv}, \rtl{pending\_mem\_req\_queue.sv}, \rtl{pending\_fp\_ops\_queue.sv}
& \ding{51}\ Fully supported \\
\addlinespace

Chapter 10. Memory Subsystem &
\rtl{mem\_unit.sv}, \rtl{vagu.sv}, \rtl{load\_store\_queue.sv}, \rtl{store\_buffer.sv}, \rtl{pending\_mem\_req\_queue.sv}
& \ding{51}\ Mostly supported \\
\addlinespace

Chapter 11. Writeback Stage &
\rtl{wb\_stage/rtl/graduation\_list.sv}, completion signals coming from execution units
& $\triangle$ Partially supported \\
\addlinespace

Chapter 12. Commit Stage &
\rtl{wb\_stage/rtl/graduation\_list.sv}, \rtl{ir\_stage/rtl/rename\_table.sv}, \rtl{free\_list.sv}, scoreboard modules
& $\triangle$ Commit logic exists but there is no standalone ROB module \\
\addlinespace

Chapter 13. Recovery Mechanisms &
\rtl{branch\_unit.sv}, \rtl{branch\_predictor.sv}, \rtl{return\_address\_stack.sv}, \rtl{control\_unit.sv}, rename/free-list recovery logic
& $\triangle$ Supported through distributed logic rather than one recovery module \\

\bottomrule
\end{longtable}
\endgroup

\vspace{1cm}

\noindent\textbf{Legend:}
\begin{itemize}[leftmargin=*]
  \item[\ding{51}] Complete RTL coverage located and matched to the documentation.
  \item[$\triangle$] Functionality exists but is distributed across several modules rather than a single dedicated file.
\end{itemize}

% ============================================================
\chapter{Abbreviations and Acronyms}

\begingroup
\renewcommand{\arraystretch}{1.2}
\begin{longtable}{@{}p{0.15\textwidth}p{0.27\textwidth}p{0.50\textwidth}@{}}
\toprule
\textbf{Abbr.} & \textbf{Full Form} & \textbf{Description} \\
\midrule
\endhead
ALU & Arithmetic Logic Unit & Executes integer arithmetic and logical operations. \\
AGU & Address Generation Unit & Computes effective addresses for memory instructions. \\
BTB & Branch Target Buffer & Stores predicted branch target addresses. \\
BHT & Branch History Table & Stores branch prediction history. \\
CSR & Control and Status Register & Special-purpose processor registers used for control and exception handling. \\
CFI & Control-Flow Integrity & Security mechanism preventing illegal control-flow transfers. \\
CPI & Cycles Per Instruction & Average execution cycles required per instruction. \\
EX & Execute Stage & Pipeline stage where arithmetic, logical, and branch operations are executed. \\
FP & Floating Point & Arithmetic performed on floating-point numbers. \\
FPU & Floating Point Unit & Dedicated execution unit for floating-point instructions. \\
FSM & Finite State Machine & Hardware state machine controlling logic behavior. \\
GHR & Global History Register & Stores global branch history for branch prediction. \\
IF & Instruction Fetch & Pipeline stage responsible for fetching instructions. \\
ID & Instruction Decode & Pipeline stage where instructions are decoded. \\
IQ & Instruction Queue & Buffer storing renamed instructions awaiting execution. \\
IR & Instruction Rename & Pipeline stage performing register renaming. \\
ISA & Instruction Set Architecture & Software-visible processor architecture. \\
LSQ & Load Store Queue & Tracks pending memory operations while preserving memory ordering. \\
LPAD & Landing Pad Instruction & Zicfilp instruction used to validate indirect branch targets. \\
MEM & Memory Stage & Pipeline stage responsible for memory accesses. \\
OoO & Out-of-Order & Execution model allowing instructions to execute before older instructions when dependencies permit. \\
PC & Program Counter & Address of the current instruction. \\
PE & Processing Element & Basic computational element, typically used in vector or matrix accelerators. \\
RAT & Register Alias Table & Maps architectural registers to physical registers during renaming. \\
ROB & Reorder Buffer & Tracks in-flight instructions and commits them in program order. \\
RR & Register Read & Pipeline stage where source operands are read. \\
RTL & Register Transfer Level & Hardware description abstraction used in Verilog/SystemVerilog. \\
RAS & Return Address Stack & Hardware stack predicting return instruction targets. \\
SIMD & Single Instruction Multiple Data & Parallel execution model processing multiple data elements simultaneously. \\
SQ & Store Queue & Queue buffering store operations before retirement. \\
SS & Shadow Stack & Protected stack used by Zicfiss to validate return addresses. \\
SV & SystemVerilog & Hardware Description Language used for RTL implementation. \\
TB & Testbench & Verification environment for RTL simulation. \\
WB & Writeback & Pipeline stage where execution results are written back. \\
XLEN & Integer Register Width & Width of the architectural integer registers (e.g., 32-bit or 64-bit). \\
\bottomrule
\end{longtable}
\endgroup

% ============================================================
\chapter{Pipeline Stage Abbreviations}

\begin{center}
\renewcommand{\arraystretch}{1.3}
\begin{tabular}{@{}ll@{}}
\toprule
\textbf{Stage} & \textbf{Meaning} \\
\midrule
IF1 & Instruction Fetch Stage 1 \\
IF2 & Instruction Fetch Stage 2 \\
ID & Instruction Decode \\
IR & Instruction Rename \\
RR & Register Read \\
EX & Execute \\
MEM & Memory Access \\
WB & Writeback \\
Commit & Instruction Retirement / Graduation \\
\bottomrule
\end{tabular}
\end{center}

% ============================================================
\chapter{RTL Module Abbreviations}

\begingroup
\renewcommand{\arraystretch}{1.2}
\begin{longtable}{@{}p{0.35\textwidth}p{0.55\textwidth}@{}}
\toprule
\textbf{Module} & \textbf{Description} \\
\midrule
\endhead
\rtl{datapath.sv} & Top-level datapath implementation \\
\rtl{control\_unit.sv} & Processor control logic \\
\rtl{if\_stage\_1.sv} & First instruction fetch stage \\
\rtl{if\_stage\_2.sv} & Second instruction fetch stage \\
\rtl{decoder.sv} & Instruction decoder \\
\rtl{immediate.sv} & Immediate value generator \\
\rtl{rename\_table.sv} & Register alias table \\
\rtl{free\_list.sv} & Physical register allocator \\
\rtl{instruction\_queue.sv} & Issue queue \\
\rtl{regfile.sv} & Physical register file \\
\rtl{exe\_stage.sv} & Top-level execution stage \\
\rtl{branch\_unit.sv} & Branch execution logic \\
\rtl{alu.sv} & Arithmetic Logic Unit \\
\rtl{div\_unit.sv} & Integer divider \\
\rtl{mul\_unit.sv} & Integer multiplier \\
\rtl{mem\_unit.sv} & Memory execution unit \\
\rtl{load\_store\_queue.sv} & Load/Store Queue \\
\rtl{pending\_mem\_req\_queue.sv} & Pending memory request tracker \\
\rtl{pending\_fp\_ops\_queue.sv} & Pending floating-point operations \\
\rtl{score\_board\_scalar.sv} & Scalar dependency scoreboard \\
\rtl{score\_board\_simd.sv} & SIMD dependency scoreboard \\
\rtl{graduation\_list.sv} & Instruction commit list \\
\rtl{csr\_interface.sv} & CSR interface logic \\
\bottomrule
\end{longtable}
\endgroup

% ============================================================
\chapter{RISC-V Extensions Referenced}

\begin{center}
\renewcommand{\arraystretch}{1.3}
\begin{tabular}{@{}ll@{}}
\toprule
\textbf{Extension} & \textbf{Description} \\
\midrule
RV64 & 64-bit RISC-V base ISA \\
M & Integer Multiply/Divide Extension \\
A & Atomic Extension \\
F & Single-Precision Floating Point \\
D & Double-Precision Floating Point \\
C & Compressed Instructions \\
V & Vector Extension \\
Zfa & Additional Floating-Point Instructions \\
Zfh & Half-Precision Floating Point \\
Zvfhmin & Minimal Half-Precision Vector Extension \\
Zicfilp & Control-Flow Integrity Landing Pad Extension \\
Zicfiss & Control-Flow Integrity Shadow Stack Extension \\
\bottomrule
\end{tabular}
\end{center}

% ============================================================
\chapter{Execution Unit Components}

\begin{center}
\renewcommand{\arraystretch}{1.3}
\begin{tabular}{@{}ll@{}}
\toprule
\textbf{Component} & \textbf{Function} \\
\midrule
ALU & Integer arithmetic and logic \\
Branch Unit & Branch comparison and target computation \\
Divider & Integer division operations \\
Multiplier & Integer multiplication operations \\
Memory Unit & Load and store execution \\
SIMD Unit & Vector arithmetic execution \\
FPU & Floating-point execution \\
AGU & Effective address generation \\
Scoreboard & Dependency tracking \\
Pending Queues & Track unfinished operations \\
\bottomrule
\end{tabular}
\end{center}
